\section{Иллюстративные случаи расхождения}
\label{sec:heldout-cases-ru}

Эти четыре примера отобраны после генерации по зафиксированному детерминированному правилу: для каждого контраста меток выбиралась задача с наименьшим числовым идентификатором в каждом направлении расхождения, если обе программы были получены и прошли критерии допуска. Примеры являются описательными и не входят в четыре основных теста. Они не оценивают частоты расхождений и не позволяют устанавливать ненаблюдаемый процесс генерации.

Для \texttt{BigCodeBench/428} ролевая программа в режиме одного вызова прошла тесты, а нейтральная нет. Обе программы выполняли требуемое внешнее объединение и масштабирование числовых признаков из \texttt{df1}. Ролевая программа сначала масштабировала копию \texttt{df1}, а затем выполняла одно объединение, сохраняя сначала столбцы левого датафрейма. Нейтральная программа сначала выполняла объединение, а затем удаляла и вставляла масштабированные столбцы вторым объединением; результат имел порядок \texttt{id, feature4, feature5, feature1, feature2, feature3}. В тесте \texttt{test\_case\_1} требовался порядок \texttt{id, feature1, feature2, feature3, feature4, feature5}, поэтому возникло расхождение списков. В условии описаны операции, но порядок столбцов прямо не задан; этот конкретный контракт задаётся тестом.

Для \texttt{BigCodeBench/71} нейтральная программа в режиме одного вызова прошла тесты, а ролевая нет. Ролевая реализация использовала выборочное стандартное отклонение, \texttt{Series.std(ddof=1)}, тогда как нейтральная реализация и эталонное решение использовали генеральное стандартное отклонение, \texttt{np.std} с \texttt{ddof=0}. Ролевая программа завершилась ошибками в \texttt{test\_case\_2}, \texttt{test\_case\_3}, \texttt{test\_case\_4} и \texttt{test\_case\_5}; зафиксированные расхождения составили соответственно $11.900380145988935$ и $11.017611134090766$, $8.568005268772557$ и $8.01463505095522$, NaN и $0.0$, а также $47.95684491664328$ и $46.07544623291379$. В тексте задания не уточнена степень свободы, тогда как исполняемый контракт однозначно требует генеральное стандартное отклонение.

Для \texttt{BigCodeBench/167} ролевая программа в режиме трёх вызовов прошла тесты, а нейтральная нет. Нейтральная реализация выбирала \texttt{max(2,min(5,num\_types))} сегментов. Поэтому при \texttt{num\_types=10} она создавала 50 патчей, а при \texttt{num\_types=1} --- 2; ролевая реализация использовала ровно \texttt{num\_types} сегментов и создавала 100 и 1 патч соответственно. В отчёте штатного оценщика указаны ошибки \texttt{test\_case\_3}, $50\mathrel{\ne}100$, и \texttt{test\_case\_4}, $2\mathrel{\ne}1$. Условие требует категории и сегментированные значения, но не задаёт число сегментов явно; тест требует равенства числа сегментов \texttt{num\_types}.

Для \texttt{BigCodeBench/31} нейтральная программа в режиме трёх вызовов прошла тесты, а ролевая нет. Нейтральная реализация сохраняла порядок первого появления элементов в \texttt{Counter}; ролевая сортировала элементы сначала по убыванию частоты, а затем лексикографически. В тестовой строке \texttt{\$is} встречается три раза, \texttt{\$second} --- один раз, а \texttt{\$sentence} --- два раза; тест ожидает значения $3,1,2$ в порядке первого появления категорий. Сортировка перемещала \texttt{\$sentence} перед \texttt{\$second}, и в \texttt{test\_case\_2} было получено сообщение ``Expected value '1.0', but got '2'.'' Условие требует график частот, но не задаёт порядок категорий, поэтому это расхождение исполняемого порядка, а не свидетельство причинного эффекта ролей.

Описания основаны только на запросах, программах и результатах штатного оценщика, приведённых для этих случаев. Они фиксируют наблюдаемое поведение кода и не делают выводов о скрытом рассуждении, независимых внутренних агентах или репрезентативных частотах в совокупности задач.
